\documentclass[11pt]{article}

\usepackage[a4paper,margin=2.5cm]{geometry}
\usepackage{amsmath,amssymb,mathtools}
\usepackage[hidelinks]{hyperref}
\usepackage[nameinlink,noabbrev]{cleveref}

\newcommand{\dd}{\mathrm{d}}
\newcommand{\Ham}{\mathcal{H}}
\newcommand{\Lag}{\mathcal{L}}
\newcommand{\RBL}{R}

\title{Timelike Geodesics in ADM Variables}
\author{}
\date{}

\begin{document}
\maketitle

\begin{abstract}
These notes derive a first-order Hamiltonian description of a timelike
geodesic using the lapse, shift, and spatial metric of a (3+1)
decomposition.  The resulting equations require precisely the ADM fields
\(\alpha,\beta^i,\gamma_{ij}\) and their spatial derivatives.  They therefore
provide a natural mathematical formulation for comparing exact, Hermite, and
Lagrange-interpolated metric data along a particle trajectory.  Unit rest mass
and signature \((-+++)\) are used throughout.
\end{abstract}

\tableofcontents

\section{Geodesics and canonical momentum}

Let \(x^\mu(\tau)\) be a future-directed timelike curve parametrized by proper
time.  Its four-velocity \(u^\mu=\dd x^\mu/\dd\tau\) satisfies
\begin{equation}
 g_{\mu\nu}u^\mu u^\nu=-1.
 \label{eq:normalization}
\end{equation}
Extremizing proper time gives the geodesic equation
\begin{equation}
 \frac{\dd u^\mu}{\dd\tau}
 +\Gamma^\mu{}_{\rho\sigma}u^\rho u^\sigma=0,
 \qquad
 \Gamma^\mu{}_{\rho\sigma}
 =\frac12g^{\mu\nu}
 \left(\partial_\rho g_{\nu\sigma}
      +\partial_\sigma g_{\nu\rho}
      -\partial_\nu g_{\rho\sigma}\right).
 \label{eq:christoffel-geodesic}
\end{equation}
Thus a geodesic probes both the metric and its first derivatives.

An equivalent canonical description uses the covariant momentum
\(p_\mu=u_\mu\) and the four-dimensional Hamiltonian
\begin{equation}
 \Ham_4(x,p)=\frac12g^{\mu\nu}p_\mu p_\nu.
 \label{eq:four-hamiltonian}
\end{equation}
Hamilton's equations with affine parameter \(\tau\) are
\begin{equation}
 \frac{\dd x^\mu}{\dd\tau}=g^{\mu\nu}p_\nu,
 \qquad
 \frac{\dd p_\mu}{\dd\tau}
 =-\frac12(\partial_\mu g^{\rho\sigma})p_\rho p_\sigma,
 \label{eq:four-hamilton-equations}
\end{equation}
subject to the mass-shell constraint
\begin{equation}
 g^{\mu\nu}p_\mu p_\nu=-1,
 \qquad \Ham_4=-\frac12.
 \label{eq:mass-shell}
\end{equation}

\section{The ADM decomposition}

Write the line element as
\begin{equation}
 \dd s^2=-\alpha^2\dd t^2
 +\gamma_{ij}(\dd x^i+\beta^i\dd t)
                  (\dd x^j+\beta^j\dd t).
 \label{eq:adm-line-element}
\end{equation}
Here \(\alpha\) is the lapse, \(\beta^i\) the contravariant shift, and
\(\gamma_{ij}\) the positive-definite spatial metric.  The spacetime metric
and its inverse are
\begin{equation}
 g_{\mu\nu}=
 \begin{pmatrix}
 -\alpha^2+\beta_k\beta^k & \beta_j\\
 \beta_i & \gamma_{ij}
 \end{pmatrix},
 \qquad
 g^{\mu\nu}=
 \begin{pmatrix}
 -\alpha^{-2} & \alpha^{-2}\beta^j\\
 \alpha^{-2}\beta^i &
 \gamma^{ij}-\alpha^{-2}\beta^i\beta^j
 \end{pmatrix},
 \label{eq:adm-matrices}
\end{equation}
where \(\beta_i=\gamma_{ij}\beta^j\) and
\(\gamma^{ik}\gamma_{kj}=\delta^i{}_j\).  In deriving the coordinate-time
system below, \(\alpha>0\) and \(t\) is assumed to be a valid time coordinate
on the region traversed by the geodesic.

Substitution into the mass-shell condition gives
\begin{equation}
 \gamma^{ij}p_i p_j
 -\frac{(p_t-\beta^i p_i)^2}{\alpha^2}=-1.
 \label{eq:adm-mass-shell}
\end{equation}
Define
\begin{equation}
 W=\sqrt{1+\gamma^{ij}p_i p_j}.
 \label{eq:W-definition}
\end{equation}
Future-directed motion has \(u^t>0\), which selects
\begin{equation}
 p_t-\beta^i p_i=-\alpha W.
 \label{eq:future-branch}
\end{equation}
Solving for \(-p_t\) produces the coordinate-time Hamiltonian
\begin{equation}
 \boxed{\Ham(t,x^i,p_i)=-p_t=\alpha W-\beta^i p_i.}
 \label{eq:reduced-hamiltonian}
\end{equation}
Indeed, restricting the canonical action to the mass shell and using \(t\)
as parameter gives
\begin{equation}
 \int p_\mu\,\dd x^\mu
 =\int\left(p_i\frac{\dd x^i}{\dd t}+p_t\right)\dd t
 =\int\left(p_i\frac{\dd x^i}{\dd t}-\Ham\right)\dd t,
 \label{eq:reduced-action}
\end{equation}
which is the standard first-order action with Hamiltonian \(-p_t\).
The three canonical momenta are covariant components \(p_i\), not coordinate
velocities or contravariant momenta.

\section{Coordinate-time evolution equations}

On any region where \(t\) is a valid time coordinate, Hamilton's equations for
\cref{eq:reduced-hamiltonian} are
\begin{equation}
 \boxed{
 \frac{\dd x^i}{\dd t}
 =\frac{\alpha\gamma^{ij}p_j}{W}-\beta^i
 }
 \label{eq:position-equation}
\end{equation}
and
\begin{equation}
 \boxed{
 \frac{\dd p_i}{\dd t}
 =-W\,\partial_i\alpha
  -\frac{\alpha}{2W}(\partial_i\gamma^{jk})p_jp_k
  +(\partial_i\beta^j)p_j.
 }
 \label{eq:momentum-equation-inverse}
\end{equation}
These equations remain Hamiltonian if the ADM fields depend explicitly on
time, although then \(\Ham\) need not be conserved.

It is often preferable to differentiate the supplied covariant spatial
metric rather than its inverse.  Differentiating
\(\gamma^{jm}\gamma_{mk}=\delta^j{}_k\) gives
\begin{equation}
 \partial_i\gamma^{jk}
 =-\gamma^{jm}(\partial_i\gamma_{mn})\gamma^{nk}.
 \label{eq:inverse-derivative}
\end{equation}
With \(q^i=\gamma^{ij}p_j\), the momentum equation becomes
\begin{equation}
 \boxed{
 \frac{\dd p_i}{\dd t}
 =-W\,\partial_i\alpha
  +\frac{\alpha}{2W}q^mq^n\partial_i\gamma_{mn}
  +p_j\partial_i\beta^j.
 }
 \label{eq:momentum-equation-covariant}
\end{equation}
Consequently the geodesic right-hand side requires exactly
\begin{equation}
 \alpha,\quad\beta^i,\quad\gamma_{ij},
 \qquad
 \partial_k\alpha,\quad\partial_k\beta^i,
 \quad\partial_k\gamma_{ij}.
 \label{eq:required-data}
\end{equation}
No Christoffel symbols need to be formed.  For a time-dependent metric the
explicit time dependence of these fields enters \(\Ham(t,x,p)\); canonical
evolution of \(x^i,p_i\) still uses their spatial derivatives.

The relation to proper time follows directly from
\cref{eq:future-branch}:
\begin{equation}
 W=\alpha u^t,
 \qquad
 \frac{\dd\tau}{\dd t}=\frac{\alpha}{W}.
 \label{eq:proper-coordinate-time}
\end{equation}

\section{Initial data}

Canonical initial data consist of \(x^i\) and \(p_i\).  If a coordinate
three-velocity \(v^i=\dd x^i/\dd t\) is specified instead, normalization of
the four-velocity yields
\begin{equation}
 u^t=
 \left[\alpha^2-
 \gamma_{ij}(v^i+\beta^i)(v^j+\beta^j)\right]^{-1/2},
 \label{eq:ut-from-velocity}
\end{equation}
provided the bracket is positive.  The canonical momentum is then
\begin{equation}
 p_i=u_i=u^t\gamma_{ij}(v^j+\beta^j).
 \label{eq:momentum-from-velocity}
\end{equation}
Conversely,
\cref{eq:position-equation} recovers \(v^i\) from \(p_i\).

\section{Conserved quantities in a stationary, axisymmetric metric}

If the metric is stationary, \(\partial_t g_{\mu\nu}=0\), then
\begin{equation}
 E=-p_t=\Ham
 \label{eq:energy}
\end{equation}
is conserved.  If it is also axisymmetric, \(p_\phi=L_z\) is conserved.  In
Cartesian coordinates adapted to rotations about the \(z\)-axis,
\begin{equation}
 L_z=x p_y-y p_x.
 \label{eq:angular-momentum}
\end{equation}
Kerr geodesics additionally possess the Carter constant.  In
Boyer--Lindquist coordinates, for unit rest mass, one common convention is
\begin{equation}
 Q=p_\theta^2+\cos^2\theta
 \left[a^2(1-E^2)+\frac{L_z^2}{\sin^2\theta}\right].
 \label{eq:carter}
\end{equation}
The quantities \(E\), \(L_z\), and \(Q\) provide independent geometric
diagnostics of a computed trajectory.  A mass-shell residual is also useful
if it is formed with an independently evolved or conserved \(p_t\).  It
vanishes algebraically if one simply resets \(p_t=-\Ham(x,p)\) at every point.
Conservation of an interpolated Hamiltonian should be distinguished from
agreement with the corresponding exact invariant.

\section{Equatorial circular Kerr data}

For Kerr mass \(M\), spin \(a>0\), and a prograde timelike circular orbit of
Boyer--Lindquist radius \(\RBL_0\), the angular frequency measured in the
coordinate \(t\) is
\begin{equation}
 \Omega=\frac{\sqrt{M}}{\RBL_0^{3/2}+a\sqrt{M}}.
 \label{eq:circular-frequency}
\end{equation}
The radius must lie in the region where the circular timelike solution exists;
choosing it outside the prograde innermost stable circular orbit gives a stable
reference orbit.
At the equator the four-velocity has the form
\begin{equation}
 u^\mu=u^t(1,0,0,\Omega),
 \qquad
 u^t=\left[-g_{tt}-2g_{t\phi}\Omega
                 -g_{\phi\phi}\Omega^2\right]^{-1/2}.
 \label{eq:circular-four-velocity}
\end{equation}

For the modified quasi-isotropic radial coordinate used in the accompanying
Kerr notes,
\begin{equation}
 \RBL=\frac{(r+s)^2}{r},
 \qquad s=\frac{r_+}{4},
 \qquad r_+=M+\sqrt{M^2-a^2}.
 \label{eq:les-map}
\end{equation}
The exterior branch corresponding to \(\RBL_0>r_+=4s\) is
\begin{equation}
 r_0=\frac{\RBL_0-2s+\sqrt{\RBL_0(\RBL_0-4s)}}{2}.
 \label{eq:les-inverse}
\end{equation}
At the Cartesian point \((x,y,z)=(r_0,0,0)\), the circular coordinate velocity
is
\begin{equation}
 (v^x,v^y,v^z)=(0,r_0\Omega,0).
 \label{eq:circular-cartesian-velocity}
\end{equation}
Equations~\eqref{eq:ut-from-velocity} and
\eqref{eq:momentum-from-velocity} then determine the canonical initial
momentum without requiring a separate transformation formula for \(p_i\).

\section{Why interpolation smoothness matters}

A metric reconstructed independently in each grid cell defines a piecewise
Hamiltonian.  If neighboring reconstructions agree only in value, the
Hamiltonian is continuous but its spatial derivatives can jump across a cell
boundary.  The force in \cref{eq:momentum-equation-covariant} can therefore
jump as a trajectory crosses that boundary.  Repeated crossings may turn
small local defects into phase error or drift in the exact invariants.

A quintic Hermite reconstruction prescribes the value, first derivative, and
second derivative at both ends of every one-dimensional cell.  When adjacent
cells share the same endpoint data, the piecewise field is \(C^2\).  The
Hamiltonian force constructed from it is then at least continuous.  A local
Lagrange polynomial is smooth within a fixed stencil, but changing the stencil
at a grid boundary generally does not match derivatives of the old and new
polynomials.  This distinction---global smoothness of the piecewise
reconstruction rather than smoothness of each individual polynomial---is the
central reason to study trajectories rather than isolated interpolation
points.

Three mathematically distinct trajectories can therefore be compared:
\begin{enumerate}
 \item the exact trajectory generated by analytic ADM fields and derivatives;
 \item the trajectory generated by a globally \(C^2\), piecewise Hermite
       reconstruction of those fields and its derivatives;
 \item the trajectory generated by locally Lagrange-interpolated fields and a
       chosen derivative construction.
\end{enumerate}
Comparisons of \(x^i(t)\), \(p_i(t)\), the instantaneous Hamiltonian vector
field, and the exact invariants separate local force defects from their
accumulated dynamical consequences.

\section{Coordinate limitation at the Kerr throat}

Assigning the signed lapse
\begin{equation}
 \alpha=(r-s)\sqrt{\frac{Q\Sigma}{rA}}
 \label{eq:signed-lapse}
\end{equation}
makes the lapse smooth on the doubled spatial slice and removes the artificial
cusp produced by \(\lvert r-s\rvert\).  It does not turn the modified
quasi-isotropic coordinates into a horizon-penetrating spacetime chart.  At
\(r=s\), the lapse
vanishes and the four-metric determinant vanishes in these coordinates.
Therefore a geodesic test formulated with coordinate time \(t\) should remain
in the regular exterior, where \(\alpha>0\).  The negative lapse on the inner
sheet supplies a smooth doubled-slice extension for interpolation; it should
not be inserted into the future-branch argument of
\cref{eq:future-branch} as though it were a positive ADM lapse in one connected
spacetime chart.  Interpolation stencils may extend across the throat when
sampling these smooth fields, but a physical trajectory cannot be evolved
through the horizon using this coordinate chart.

\end{document}
